\begin{table}[htbp]
\centering
\caption{Execution-environment provenance for the preserved V1.0 runs and V1.1 development.}
\label{tab:environment-comparison}
\small
\begingroup
\setlength{\tabcolsep}{5pt}
\renewcommand{\arraystretch}{1.2}
\begin{tabularx}{\linewidth}{@{}>{\raggedright\arraybackslash}p{0.17\linewidth}>{\raggedright\arraybackslash}X>{\raggedright\arraybackslash}X>{\raggedright\arraybackslash}X@{}}
\toprule
Research state & V1.0 initial run & V1.0 confirmatory repetitions & V1.1 development environment \\
\midrule
Operating system & Windows 11 build 26200 & Debian GNU/Linux, kernel 6.12.101 & Windows 10 Pro 22H2 build 19045 \\
CPU & Intel Core i5-1335U & Intel Core i5-1335U & AMD Ryzen 5 5600 \\
Topology & 12 logical processors & 10 cores / 12 logical processors & 6 cores / 12 logical processors \\
RAM (GiB) & 15.7 & 15.4 & 31.9 \\
Inference device & CPU; Intel Iris Xe present & CPU & NVIDIA RTX 3060, 12 GiB VRAM \\
\bottomrule
\end{tabularx}
\endgroup
\par\vspace{2pt}{\raggedright\footnotesize\textit{Note:} The V1.1 row is an environment inventory, not evidence that a V1.1 experiment had already run. Cross-machine latency is descriptive only because hardware, operating system, driver, runtime version, and CPU/GPU placement differ.}
\par\vspace{2pt}{\raggedright\fontsize{9}{11}\selectfont\textbf{Source:} Research data (2026).}
\end{table}
